\section{\texorpdfstring{Derived Categories:\allowbreak{} full subcategories and generators}{Derived  full subcategories and generators}}\label{proofs-10691-11108-derived-categories-full-subcategories-and-generators}

Private correction review,\allowbreak{} 2026-09-30. Authority:\allowbreak{} The Stacks Project,\allowbreak{} commit a04446e5\allowbreak{}7ec1fbc2\allowbreak{}52a871af\allowbreak{}cec7752f\allowbreak{}b2807b14, derived.tex 10691–\allowbreak{}11108,\allowbreak{} with its exact-functor definition at 190–\allowbreak{}203. The receiver equiv.tex is frozen separately in authority-equiv-part18.tex and current-equiv-part18.tex. Those two files are identical. Source identities,\allowbreak{} original reports,\allowbreak{} exact replacements and actual reading coverage are bound by PART18\_\allowbreak{}CHECKPOINT.json and its associated ledgers.

The original objects,\allowbreak{} shifts,\allowbreak{} extension triangles and tensor factors are retained below. Sections 6 and 7 are editorial consequences,\allowbreak{} not replacement source statements or novelty claims. The corrections in Sections 2 and 8 are distinct from the seven earlier corrections reconciled here.

\subsection{\texorpdfstring{1. The operations,\allowbreak{} their exact witnesses and finite closure}{1. The  their exact witnesses and finite closure}}\label{proofs-10691-11108-the-operations-their-exact-witnesses-and-finite-closure}

Write T for the given triangulated category. A full subcategory A is an actual subset of Ob(T)\allowbreak{},\allowbreak{} as specified at original 10695–\allowbreak{}10696. Thus A[a,\allowbreak{}b]\allowbreak{} contains precisely the objects A[-i]\allowbreak{} for A in A and integer a <\allowbreak{}= i <\allowbreak{}= b; it is not implicitly closed under arbitrary isomorphism. In contrast,\allowbreak{} smd(A)\allowbreak{} and add(A)\allowbreak{} include objects isomorphic to actual direct summands and finite direct sums,\allowbreak{} respectively. The empty finite sum is allowed. A star B consists of the middle objects of distinguished triangles A \allowbreak{}-> X \allowbreak{}-> B \allowbreak{}-> A[1]\allowbreak{}.

For associativity,\allowbreak{} start with triangles

\[
 A\xrightarrow f X\xrightarrow p B\xrightarrow h A[1],
 \qquad X\xrightarrow g Y\xrightarrow q C\xrightarrow k X[1].
\]

Complete gf to A \allowbreak{}-> Y \allowbreak{}-> Z \allowbreak{}-> A[1]\allowbreak{}. TR4 gives a distinguished triangle B \allowbreak{}-> Z \allowbreak{}-> C \allowbreak{}-> B[1]\allowbreak{},\allowbreak{} whose last map is p[1]\allowbreak{}k,\allowbreak{} and maps of triangles having first two components (id\_\allowbreak{}A,\allowbreak{}g)\allowbreak{} and (f,\allowbreak{}id\_\allowbreak{}Y)\allowbreak{}. Therefore Y belongs to A star (B star C)\allowbreak{}.

Conversely take triangles

\[
 A\xrightarrow f Y\xrightarrow g Z\xrightarrow h A[1],
 \qquad B\xrightarrow u Z\xrightarrow v C\xrightarrow w B[1].
\]

Apply TR4 to Y --g-\allowbreak{}-> Z --v-\allowbreak{}-> C,\allowbreak{} using the rotated triangles (Y,\allowbreak{}Z,\allowbreak{}A[1]\allowbreak{},\allowbreak{}g,\allowbreak{}h,\allowbreak{}-f[1]\allowbreak{})\allowbreak{} and (Z,\allowbreak{}C,\allowbreak{}B[1]\allowbreak{},\allowbreak{}v,\allowbreak{}w,\allowbreak{}-u[1]\allowbreak{})\allowbreak{}. Choose a triangle (Y,\allowbreak{}C,\allowbreak{}W,\allowbreak{}vg,\allowbreak{}p,\allowbreak{}d)\allowbreak{}. Its cone triangle is

\[
 A[1]\xrightarrow\alpha W\xrightarrow\beta B[1]
       \xrightarrow{-(hu)[1]} A[2].
\]

Put X\allowbreak{}=W[-1]\allowbreak{}. Inverse rotation gives X --(-d[-1]\allowbreak{})\allowbreak{}-\allowbreak{}-> Y --vg-\allowbreak{}-> C --p-\allowbreak{}-> X[1]\allowbreak{}. Inverse shift of the cone triangle,\allowbreak{} with its triangle boundary sign,\allowbreak{} gives A --alpha[-1]\allowbreak{}-\allowbreak{}-> X --beta[-1]\allowbreak{}-\allowbreak{}-> B --hu-\allowbreak{}-> A[1]\allowbreak{}. Here maps are transported along the chosen inverse-shift identifications; no literal equality of composed shift functors is assumed. This proves the other inclusion with the original A,\allowbreak{} B,\allowbreak{} C and Y. It supplies the direction left implicit at 10733–\allowbreak{}10735.

Next let A1 \allowbreak{}-> X \allowbreak{}-> B1 \allowbreak{}-> A1[1]\allowbreak{},\allowbreak{} with maps f,\allowbreak{}g,\allowbreak{}h,\allowbreak{} be distinguished. Choose the actual complement witnesses for membership in smd(A)\allowbreak{} and smd(B)\allowbreak{}:\allowbreak{} isomorphisms A1\allowbreak{}+ A2 \allowbreak{}-> A0 in A and B1\allowbreak{}+ B2 \allowbreak{}-> B0 in B. The sum of this triangle with A2 --id-\allowbreak{}-> A2 \allowbreak{}-> 0 and 0 \allowbreak{}-> B2 --id-\allowbreak{}-> B2 gives,\allowbreak{} after permuting the middle summands,\allowbreak{} the distinguished triangle

\[
 A_1\oplus A_2
 \xrightarrow{(a_1,a_2)\mapsto(a_2,f a_1,0)}
 A_2\oplus X\oplus B_2
 \xrightarrow{(a_2,x,b_2)\mapsto(gx,b_2)}
 B_1\oplus B_2
 \xrightarrow{(b_1,b_2)\mapsto(hb_1,0)}
 (A_1\oplus A_2)[1].
\]

Transport its endpoints to A0,\allowbreak{}B0. Its middle object has X as a specified summand,\allowbreak{} proving smd(A)\allowbreak{} star smd(B)\allowbreak{} subset smd(A star B)\allowbreak{}. Actual complement witnesses compose:\allowbreak{} if Y\allowbreak{}=X\allowbreak{}+X' and Z\allowbreak{}=Y\allowbreak{}+Y',\allowbreak{} then Z\allowbreak{}=X\allowbreak{}+(X'\allowbreak{}+Y')\allowbreak{}. Hence smd(smd(A)\allowbreak{})\allowbreak{}\allowbreak{}=smd(A)\allowbreak{}. The reverse inclusion in the asserted equality after another smd follows from A subset smd(A)\allowbreak{} and B subset smd(B)\allowbreak{}. No assumption that every idempotent in T splits has entered this proof.

Finite sums of distinguished triangles are distinguished. Combining the finite lists of summands proves add(A)\allowbreak{} star add(B)\allowbreak{} is closed under finite sums; combining complement witnesses proves the same for smd(add(A)\allowbreak{})\allowbreak{}. The assertions about direct sums at 10763 and 10783 are used in this finite sense,\allowbreak{} as in their displayed binary-sum proof. They do not establish closure under infinite coproducts. For example,\allowbreak{} every object in smd(add(\{Z\})\allowbreak{})\allowbreak{} in D(Ab)\allowbreak{} has finitely generated H\textasciicircum{}0,\allowbreak{} whereas the countable direct sum of Z does not.

Set C\_\allowbreak{}n(A)\allowbreak{}\allowbreak{}=smd(add(A)\allowbreak{}\textasciicircum{}\{star n\})\allowbreak{},\allowbreak{} n>\allowbreak{}=1. The preceding arguments prove that C\_\allowbreak{}n(A)\allowbreak{} is strictly full,\allowbreak{} closed under finite sums and actual summands. Associativity and the smd equality give

\[
 C_{n+m}(A)=smd(C_n(A)\star C_m(A)).
\]

Since 0 is an empty finite sum,\allowbreak{} adding a zero factor through X --id-\allowbreak{}-> X \allowbreak{}-> 0 proves C\_\allowbreak{}n(A)\allowbreak{} subset C\_\allowbreak{}\{n\allowbreak{}+1\}(A)\allowbreak{}. All these operations are monotone in A.

For an increasing sequence A\_\allowbreak{}j,\allowbreak{} every operation in 10823–\allowbreak{}10841 commutes with union:\allowbreak{} a shift or summand has one witness in some A\_\allowbreak{}j; a finite sum has finitely many such witnesses; a fixed n-fold extension has finitely many leaves in its extension tree. Their indices have a common upper bound. The same triangles,\allowbreak{} maps,\allowbreak{} complements and sums then witness membership using that A\_\allowbreak{}j. The reverse inclusions follow from monotonicity. This is an actual finite-witness proof and makes no assertion for infinite sums.

\subsection{2. Exact functors and literal object sets --- DERIVED-NEW-037}\label{proofs-10691-11108-exact-functors-and-literal-object-sets-derived-new-037}

The source definition equips an exact functor F with natural isomorphisms xi\_\allowbreak{}X:\allowbreak{}F(X[1]\allowbreak{})\allowbreak{}\allowbreak{}->F(X)\allowbreak{}[1]\allowbreak{}. Iterating xi,\allowbreak{} its inverse and the chosen inverse-shift identifications gives,\allowbreak{} for every integer i,\allowbreak{}

\[
 \xi_X^{(-i)}:F(X[-i])\xrightarrow{\sim}F(X)[-i].
\]

Consequently F(A[a,\allowbreak{}b]\allowbreak{})\allowbreak{} and F(A)\allowbreak{}[a,\allowbreak{}b]\allowbreak{} have the same isomorphism closure. This is sufficient for every subsequent occurrence inside add,\allowbreak{} smd or C\_\allowbreak{}n. It does not give the literal object-set equality printed at 10800.

Here is a counterexample in the source\textquotesingle s actual category conventions. Let Q be the complex Z --id-\allowbreak{}-> Z in degrees 0 and 1 in D(Ab)\allowbreak{}. Its contracting homotopy has h\textasciicircum{}1\allowbreak{}=id\_\allowbreak{}Z and all other components zero,\allowbreak{} so id\_\allowbreak{}Q\allowbreak{}=0 in D(Ab)\allowbreak{}. Define F on every object to be Q and on every morphism to be the unique morphism Q\allowbreak{}->Q,\allowbreak{} which is simultaneously zero and id\_\allowbreak{}Q. This respects identities and composition and is additive. The unique map Q\allowbreak{}->Q[1]\allowbreak{} is an isomorphism,\allowbreak{} since both are zero objects in D(Ab)\allowbreak{}; it defines xi. Every image triangle is isomorphic to the zero triangle,\allowbreak{} so F is exact.

Take A to have the single object the zero complex,\allowbreak{} and a\allowbreak{}=b\allowbreak{}=1. F(A[1,\allowbreak{}1]\allowbreak{})\allowbreak{} has the sole object Q. F(A)\allowbreak{}[1,\allowbreak{}1]\allowbreak{} has the sole object Q[-1]\allowbreak{}. These are different complexes,\allowbreak{} with different nonzero terms,\allowbreak{} and thus different objects of D(Ab)\allowbreak{},\allowbreak{} although isomorphic there. The literal assertion is false. The repair states equality of isomorphism closures without changing either F or A.

The other four inclusions need no repair. An additive functor sends an actual decomposition X\allowbreak{}+Y \allowbreak{}-> A0 to a decomposition F(X)\allowbreak{}\allowbreak{}+F(Y)\allowbreak{} \allowbreak{}-> F(A0)\allowbreak{},\allowbreak{} and sends a finite sum to its canonically isomorphic finite sum of images. For an extension it sends the distinguished triangle to

\[
 F(A)\xrightarrow{F(f)}F(X)\xrightarrow{F(g)}F(B)
       \xrightarrow{\xi_A F(h)}F(A)[1].
\]

Thus its actual middle object belongs to F(A)\allowbreak{} star F(B)\allowbreak{}. Induction gives the n-fold version. At n\allowbreak{}=1 this is the identity inclusion,\allowbreak{} not an appeal to isomorphism closure of A.

MC-STK-ERR-0888 changes ``objects consists'' to ``objects consist''. It corrects the plural verb and is separate from the formula issue. The public source comparison is \href{https://stacks.math.columbia.edu/tag/0FX5}{the operations remark} and \href{https://stacks.math.columbia.edu/tag/05QK}{the exact-functor definition}. The frozen original TeX,\allowbreak{} not that mutable comparison,\allowbreak{} is the authority.

\subsection{3. Cohomology layers and actual complex terms}\label{proofs-10691-11108-cohomology-layers-and-actual-complex-terms}

Let A be an abelian category and let E be the given set of objects of A,\allowbreak{} embedded in D(A)\allowbreak{} in degree zero. Suppose H\textasciicircum{}i(K)\allowbreak{}\allowbreak{}=0 outside [a,\allowbreak{}b]\allowbreak{}. If a\allowbreak{}=b,\allowbreak{} the canonical truncation maps identify K with H\textasciicircum{}a(K)\allowbreak{}[-a]\allowbreak{}. They induce the identity on its sole cohomology object and are quasi-isomorphisms. This proves the repair in MC-STK-ERR-0889:\allowbreak{} the index is a,\allowbreak{} not an unbound i.

For b\textgreater a the canonical truncation triangle is

\[
 \tau_{\leq b-1}K\longrightarrow K\longrightarrow H^b(K)[-b]
                 \longrightarrow(\tau_{\leq b-1}K)[1].
\]

More precisely,\allowbreak{} its middle term is first tau\_\allowbreak{}\{\textbackslash{}leq b\}K,\allowbreak{} whose canonical map to K is a quasi-isomorphism under the stated upper bound,\allowbreak{} and the displayed triangle is its transport. Its third term is cohomology,\allowbreak{} not the last term of an arbitrary representative. For example,\allowbreak{} for the contractible complex M --id-\allowbreak{}-> M in degrees b-1,\allowbreak{}b,\allowbreak{} both K and its canonical lower truncation are zero in D(A)\allowbreak{},\allowbreak{} but K\textasciicircum{}b[-b]\allowbreak{}\allowbreak{}=M[-b]\allowbreak{} is nonzero when M is nonzero. It cannot be the third term of that triangle. This verifies MC-STK-ERR-0890.

If each H\textasciicircum{}i(K)\allowbreak{} belongs to smd(add(E)\allowbreak{})\allowbreak{},\allowbreak{} then H\textasciicircum{}i(K)\allowbreak{}[-i]\allowbreak{} belongs to C\_\allowbreak{}1(E[a,\allowbreak{}b]\allowbreak{})\allowbreak{},\allowbreak{} by shifting its actual finite-sum and complement witnesses. Induction on b-a,\allowbreak{} using Section 1,\allowbreak{} gives K in C\_\allowbreak{}\{b-a\allowbreak{}+1\}(E[a,\allowbreak{}b]\allowbreak{})\allowbreak{}. Membership H\textasciicircum{}i(K)\allowbreak{} in E is the special case used in part (1)\allowbreak{}. Zero intermediate layers cause no problem because 0 belongs to C\_\allowbreak{}1 regardless of whether E contains zero.

The omitted termwise proof uses a different triangle. If K\^{}bullet has actual terms zero outside [a,\allowbreak{}b]\allowbreak{},\allowbreak{} its top term is a subcomplex,\allowbreak{} and the degreewise split exact sequence of complexes is

\[
 0\longrightarrow K^b[-b]\longrightarrow K^\bullet
   \longrightarrow \sigma_{\leq b-1}K^\bullet\longrightarrow0.
\]

Here sigma is the stupid quotient truncation,\allowbreak{} with last term K\^{}\{b-1\} and outgoing differential zero; it is not canonical cohomological truncation. The connecting map of this sequence retains d\textasciicircum{}\{b-1\}:\allowbreak{}K\textasciicircum{}\{b-1\}\allowbreak{}->K\textasciicircum{}b with the shift sign of its triangle. Removing top terms successively proves parts (3)\allowbreak{} and (4)\allowbreak{},\allowbreak{} in descending term order,\allowbreak{} using K\textasciicircum{}i[-i]\allowbreak{} in C\_\allowbreak{}1(E[a,\allowbreak{}b]\allowbreak{})\allowbreak{}. The entire differential information remains in the extension maps. No direct-sum splitting of K into its terms or its cohomology is asserted.

\subsection{4. Degree support is independent of extension length}\label{proofs-10691-11108-degree-support-is-independent-of-extension-length}

For the homological functor in the source let H\textasciicircum{}q(X)\allowbreak{}\allowbreak{}=H(X[q]\allowbreak{})\allowbreak{}. Shift comparisons give H\textasciicircum{}q(E[-j]\allowbreak{})\allowbreak{}\allowbreak{}=H\textasciicircum{}\{q-j\}(E)\allowbreak{},\allowbreak{} through the specified natural identifications. Thus if E has support in [a,\allowbreak{}b]\allowbreak{},\allowbreak{} E[-j]\allowbreak{} has support in [a\allowbreak{}+j,\allowbreak{}b\allowbreak{}+j]\allowbreak{}. For -m<\allowbreak{}=j<\allowbreak{}=m the union of these intervals is contained in [a-m,\allowbreak{}b\allowbreak{}+m]\allowbreak{}.

Additivity preserves vanishing under finite sums and summands. In a distinguished triangle A\allowbreak{}->X\allowbreak{}->B\allowbreak{}->A[1]\allowbreak{},\allowbreak{} the exact sequence H\textasciicircum{}q(A)\allowbreak{}\allowbreak{}->H\textasciicircum{}q(X)\allowbreak{}\allowbreak{}->H\textasciicircum{}q(B)\allowbreak{} proves H\textasciicircum{}q(X)\allowbreak{}\allowbreak{}=0 whenever its two neighbors vanish. Induction proves that every X in C\_\allowbreak{}n(E[-m,\allowbreak{}m]\allowbreak{})\allowbreak{} has support in [a-m,\allowbreak{}b\allowbreak{}+m]\allowbreak{},\allowbreak{} for every n>\allowbreak{}=1. This verifies MC-STK-ERR-0891 without an unstated boundedness assumption on a chosen representative.

The former interval [-m\allowbreak{}+na,\allowbreak{}m\allowbreak{}+nb]\allowbreak{} is false:\allowbreak{} in D(k)\allowbreak{},\allowbreak{} take the singleton E\allowbreak{}=\{k[-2]\allowbreak{}\},\allowbreak{} a\allowbreak{}=b\allowbreak{}=2,\allowbreak{} m\allowbreak{}=1 and n\allowbreak{}=2. The split triangle E --id-\allowbreak{}-> E \allowbreak{}-> 0 and 0 in add(E[-1,\allowbreak{}1]\allowbreak{})\allowbreak{} put E in C\_\allowbreak{}2(E[-1,\allowbreak{}1]\allowbreak{})\allowbreak{}. Its H\^{}2 is k,\allowbreak{} outside the formerly claimed interval [3,\allowbreak{}5]\allowbreak{}. The argument also proves the more precise support statement:\allowbreak{} for any set S of degrees supporting all E and any set J of permitted shift indices,\allowbreak{} all finite extensions,\allowbreak{} sums and summands of E[-j]\allowbreak{} have support contained in \{s\allowbreak{}+j:\allowbreak{}s in S,\allowbreak{}j in J\}. Extension length does not multiply or fill gaps in that set.

\subsection{\texorpdfstring{5. Generator levels,\allowbreak{} minimality and detection}{5. Generator  minimality and detection}}\label{proofs-10691-11108-generator-levels-minimality-and-detection}

Let A\allowbreak{}=E[-infinity,\allowbreak{}infinity]\allowbreak{}. The source definition gives <E>\_\allowbreak{}1\allowbreak{}=smd(add(A)\allowbreak{})\allowbreak{}. Induction with Section 1 gives

\[
 \langle E\rangle_n=C_n(A)
 =smd(\langle E\rangle_1^{\star n})
 =\bigcup_{m\geq1}C_n(E[-m,m]).
\]

The last equality follows from the finite-witness union argument. In particular every individual membership has a finite shift window,\allowbreak{} even though no common window for all members is claimed. We also have <E>\_\allowbreak{}\{n\allowbreak{}+n'\}\allowbreak{}=smd(<E>\_\allowbreak{}n star <E>\_\allowbreak{}\{n'\})\allowbreak{}. MC-STK-ERR-0892 supplies the missing ``by'' in ``Denote by''. MC-STK-ERR-0893 removes the unmatched bracket in the formula for <E>\_\allowbreak{}1; the balanced formula above is the source\textquotesingle s definition.

The increasing union \textless E\textgreater{} of these levels is closed under finite sums and summands. A shift of a finite witness is again a finite witness with all its indices shifted,\allowbreak{} so it is closed under both shifts. If X is in level r and Y in level s,\allowbreak{} a cone triangle X\allowbreak{}->Y\allowbreak{}->Z\allowbreak{}->X[1]\allowbreak{} puts Z in level s\allowbreak{}+r. Thus this strictly full subcategory is triangulated and saturated.

If D\textquotesingle{} is another strictly full saturated triangulated subcategory containing E,\allowbreak{} it is stable under shifts,\allowbreak{} finite sums,\allowbreak{} actual summands and extensions. Each witness for every level therefore lies in D\textquotesingle. This proves minimality. The correct closure assertion is add(D')\allowbreak{} subset D',\allowbreak{} as in MC-STK-ERR-0894. The former add(D)\allowbreak{} subset D\textquotesingle{} is false:\allowbreak{} take D\allowbreak{}=D(k)\allowbreak{},\allowbreak{} E\allowbreak{}=0,\allowbreak{} and D\textquotesingle{} the strictly full subcategory of zero objects. It contains E but not k. The P02-E0085 explanation concerning a functor landing in D\textquotesingle{} does not describe this proof; the replacement is justified instead by this exact closure argument.

For the right orthogonal,\allowbreak{} assume Hom(E,\allowbreak{}K[i]\allowbreak{})\allowbreak{}\allowbreak{}=0 for every i. Shifting the source and target shows Hom(E[j]\allowbreak{},\allowbreak{}K)\allowbreak{}\allowbreak{}=0 for every j. Finite sums and summands preserve this vanishing. For A\allowbreak{}->X\allowbreak{}->B in an extension witness,\allowbreak{} exactness of Hom(B,\allowbreak{}K)\allowbreak{}\allowbreak{}->Hom(X,\allowbreak{}K)\allowbreak{}\allowbreak{}->Hom(A,\allowbreak{}K)\allowbreak{} proves the same vanishing for X. Induction over the levels proves Hom(X,\allowbreak{}K)\allowbreak{}\allowbreak{}=0 for all X in \textless E\textgreater. Conversely choose X\allowbreak{}=E[-i]\allowbreak{}. If E is classical,\allowbreak{} <E>\allowbreak{}=D,\allowbreak{} so this vanishing includes id\_\allowbreak{}K\allowbreak{}=0. The unique zero maps between K and a zero object are then inverse,\allowbreak{} proving K is zero. Hence classical generators are weak generators.

The source already proves the change-of-generator level bound:\allowbreak{} if E\textquotesingle{} is in <E>\_\allowbreak{}r,\allowbreak{} every shift,\allowbreak{} finite sum and summand of E\textquotesingle{} is in <E>\_\allowbreak{}r,\allowbreak{} so <E'>\_\allowbreak{}1 subset <E>\_\allowbreak{}r and <E'>\_\allowbreak{}n subset <E>\_\allowbreak{}\{rn\}. Thus a classical generator is strong whenever D has a strong generator. This is not counted as a new editorial consequence.

Finally consider the object property T in 11094–\allowbreak{}11108. From T(E)\allowbreak{} and the triangle E --id-\allowbreak{}-> E \allowbreak{}-> 0 one obtains T(0)\allowbreak{}. If X is isomorphic to Y,\allowbreak{} the triangle X\allowbreak{}->Y\allowbreak{}->0 shows T(X)\allowbreak{} implies T(Y)\allowbreak{}. Thus isomorphism invariance follows from the actual assumptions; it need not be added. The triangle X\allowbreak{}->0\allowbreak{}->X[1]\allowbreak{} and its inverse rotation give shift closure. The stated sum,\allowbreak{} triangle and summand assumptions now show that the objects satisfying T form a strictly full saturated triangulated subcategory containing E,\allowbreak{} so minimality proves the remark.

\subsection{6. Sparse and weighted generation --- DERIVED-UNDER-027}\label{proofs-10691-11108-sparse-and-weighted-generation-derived-under-027}

Keep the actual K and its canonical truncation maps. Let p1\textless...\textless pt be its nonzero cohomology degrees,\allowbreak{} a finite nonempty list. Suppose E is a full subcategory of D(A)\allowbreak{} and,\allowbreak{} for each pj,\allowbreak{}

\[
 H^{p_j}(K)\in C_{r_j}(E),\qquad r_j\geq1.
\]

Write a\allowbreak{}=p1 and b\allowbreak{}=pt. Then

\[
 K\in C_{r_1+\cdots+r_t}(E[a,b]).
\]

Indeed,\allowbreak{} shifting each witness gives H\textasciicircum{}\{pj\}(K)\allowbreak{}[-pj]\allowbreak{} in C\_\allowbreak{}\{rj\}(E[a,\allowbreak{}b]\allowbreak{})\allowbreak{}. Between two consecutive nonzero cohomology degrees,\allowbreak{} the canonical truncation maps are quasi-isomorphisms,\allowbreak{} since they induce identity on every retained cohomology object and zero-to-zero elsewhere. Compose these actual maps across each gap. The remaining tower has distinguished triangles

\[
 \tau_{\leq p_{j-1}}K\longrightarrow\tau_{\leq p_j}K
 \longrightarrow H^{p_j}(K)[-p_j]
 \longrightarrow(\tau_{\leq p_{j-1}}K)[1]
\]

for j>\allowbreak{}=2,\allowbreak{} after transport along those quasi-isomorphisms. Its first stage is H\textasciicircum{}\{p1\}(K)\allowbreak{}[-p1]\allowbreak{},\allowbreak{} and its last stage is canonically isomorphic to K. Apply the level-addition identity at every displayed triangle. This proves the stated sum of costs. When all cohomology is zero,\allowbreak{} K is isomorphic to zero and belongs to C\_\allowbreak{}1(E[a,\allowbreak{}b]\allowbreak{})\allowbreak{} for any chosen nonempty interval; no level zero is introduced into the source notation.

There is an equally exact termwise statement. For a given bounded complex with nonzero terms in degrees q1<...<qu,\allowbreak{} assume K\^{}\{qj\} in C\_\allowbreak{}\{sj\}(E)\allowbreak{}. Peel off its actual top term using the degreewise split sequence in Section 3,\allowbreak{} then repeat with the quotient. The resulting extension order is qu,\allowbreak{}...,\allowbreak{}q1,\allowbreak{} and its cost is s1\allowbreak{}+...\allowbreak{}+su in the shift window [q1,\allowbreak{}qu]\allowbreak{}. All intervening differentials are the original differentials in those triangles.

For the source hypothesis rj\allowbreak{}=1,\allowbreak{} this reduces the displayed width b-a\allowbreak{}+1 to the number of nonzero layers. Weighted costs allow different finite constructions for different layers. Neither bound is asserted to be minimal,\allowbreak{} and no splitting of the actual extension data is used. This consequence belongs in editorial material.

\subsection{7. Change of generator with the full shift window --- DERIVED-UNDER-028}\label{proofs-10691-11108-change-of-generator-with-the-full-shift-window-derived-under-028}

Let G,\allowbreak{}E be objects of T,\allowbreak{} let a<\allowbreak{}=b and c<\allowbreak{}=d be finite integers,\allowbreak{} and let r,\allowbreak{}n>\allowbreak{}=1. For the original witnesses satisfying

\[
 E\in C_r(G[c,d]),\qquad X\in C_n(E[a,b]),
\]

one has

\[
 X\in C_{rn}(G[a+c,b+d]).
\]

To prove this,\allowbreak{} fix i in [a,\allowbreak{}b]\allowbreak{}. Shift every triangle,\allowbreak{} finite sum,\allowbreak{} complement and isomorphism in the witness for E by [-i]\allowbreak{}. Its leaves become G[-j]\allowbreak{}[-i]\allowbreak{},\allowbreak{} j in [c,\allowbreak{}d]\allowbreak{}. The specified composition isomorphisms identify these with G[-(i\allowbreak{}+j)\allowbreak{}]\allowbreak{}. Therefore E[-i]\allowbreak{} belongs to C\_\allowbreak{}r(G[c\allowbreak{}+i,\allowbreak{}d\allowbreak{}+i]\allowbreak{})\allowbreak{},\allowbreak{} which is contained in C\_\allowbreak{}r(G[a\allowbreak{}+c,\allowbreak{}b\allowbreak{}+d]\allowbreak{})\allowbreak{}. This receiving category is strictly full,\allowbreak{} closed under finite sums and actual summands,\allowbreak{} so the same statement holds for every object of smd(add(E[a,\allowbreak{}b]\allowbreak{})\allowbreak{})\allowbreak{}.

For each of the n extension factors witnessing X,\allowbreak{} substitute that membership. Apply Section 1 successively to the same extension triangles and the final summand:\allowbreak{} their costs add to rn. Their common window remains [a\allowbreak{}+c,\allowbreak{}b\allowbreak{}+d]\allowbreak{}. It is not multiplied by n or r. Forgetting the finite windows gives the source\textquotesingle s already established unbounded rn rule; retaining the exact windows is the refinement recorded here. It does not change the source definitions.

\subsection{8. The tensor factor in the geometric receiver --- DERIVED-NEW-038}\label{proofs-10691-11108-the-tensor-factor-in-the-geometric-receiver-derived-new-038}

In equiv.tex 1784–\allowbreak{}1842,\allowbreak{} X is proper and smooth over the field k,\allowbreak{} E and G are the fixed finite locally free sheaves chosen from the diagonal construction,\allowbreak{} and

\[
 \mathcal O_\Delta\in C_n((E\boxtimes G)[-m,m]).
\]

For a perfect K,\allowbreak{} the exact functor

\[
 \Phi_K(M)=R{\rm pr}_{2,*}(L{\rm pr}_1^*K
                   \otimes_{\mathcal O_{X\times X}}^{\mathbf L}M)
\]

sends O\_\allowbreak{}Delta to K and E boxtimes G to

\[
 V_K\otimes_k G,\qquad
 V_K=R\Gamma(X,K\otimes_{\mathcal O_X}^{\mathbf L}E).
\]

The first identification is restriction to the diagonal followed by its inverse projection. For the second,\allowbreak{} flat base change from Spec(k)\allowbreak{} identifies Rpr\_\allowbreak{}\{2,\allowbreak{}*\}Lpr\_\allowbreak{}1\textasciicircum{}*(K tensor\^{}L E)\allowbreak{} with RΓ(X,\allowbreak{}K tensor\^{}L E)\allowbreak{} tensor\_\allowbreak{}k O\_\allowbreak{}X,\allowbreak{} and the projection formula tensors it with G. All schemes here are quasi-compact and separated,\allowbreak{} and all complexes have quasi-coherent cohomology; over the field k the base change is flat. The projection formula used is perfect.tex,\allowbreak{} lemma-cohomology-base-change,\allowbreak{} current 5051–\allowbreak{}5094. These are isomorphisms applied inside the strictly full C\_\allowbreak{}n,\allowbreak{} so the repair in Section 2 preserves the entire argument.

The complex K tensor\^{}L E is perfect:\allowbreak{} on each open with a bounded finite locally free representative for K,\allowbreak{} tensor it with the finite locally free E. Properness gives finitely many such opens,\allowbreak{} hence a bounded coherent cohomology complex globally. Applying perfect.tex,\allowbreak{} lemma-direct-image-coherent,\allowbreak{} current 2634–\allowbreak{}2660,\allowbreak{} to X\allowbreak{}->Spec(k)\allowbreak{} proves V\_\allowbreak{}K has finite-dimensional cohomology in a finite interval [a,\allowbreak{}b]\allowbreak{}. Its hypotheses hold:\allowbreak{} Spec(k)\allowbreak{} is Noetherian,\allowbreak{} X is of finite type,\allowbreak{} and every support is closed in proper X,\allowbreak{} hence proper. The bounds must therefore refer to H\textasciicircum{}i(X,\allowbreak{}K tensor\^{}L E)\allowbreak{},\allowbreak{} not H\textasciicircum{}i(X,\allowbreak{}K)\allowbreak{}.

Here is the precise use of the field. For any complex V of k-vector spaces write Z\textasciicircum{}q\allowbreak{}=ker d\^{}q and B\textasciicircum{}q\allowbreak{}=im d\^{}\{q-1\}. Choose vector-space complements Z\textasciicircum{}q\allowbreak{}=B\textasciicircum{}q\allowbreak{}+T\textasciicircum{}q and V\textasciicircum{}q\allowbreak{}=Z\textasciicircum{}q\allowbreak{}+S\textasciicircum{}q. The map d\^{}q restricts to an isomorphism S\textasciicircum{}q\allowbreak{}->B\textasciicircum{}\{q\allowbreak{}+1\}; T\^{}q maps isomorphically to H\textasciicircum{}q(V)\allowbreak{}. Thus V is the direct sum of its zero-differential T\textasciicircum{}q[-q]\allowbreak{} pieces and the two-term disks S\^{}q --d\textasciicircum{}q-\allowbreak{}-> B\textasciicircum{}\{q\allowbreak{}+1\}. Every degree has just the stated finite number of types of pieces,\allowbreak{} so this decomposition is well-defined even when V itself is unbounded. On each disk the inverse of d\^{}q is a contraction. The projection onto the T pieces and their inclusions give an explicit homotopy equivalence,\allowbreak{} with the contractions accounting for every discarded disk.

Apply this to V\_\allowbreak{}K. Its nonzero T pieces are finite-dimensional and confined to [a,\allowbreak{}b]\allowbreak{},\allowbreak{} so V\_\allowbreak{}K tensor\_\allowbreak{}k G is isomorphic to the finite sum of H\textasciicircum{}q(V\_\allowbreak{}K)\allowbreak{} tensor\_\allowbreak{}k G[-q]\allowbreak{},\allowbreak{} a member of add(G[a,\allowbreak{}b]\allowbreak{})\allowbreak{}. Tensoring preserves the displayed contraction. The same shift calculation as Section 7 gives

\[
 C_n((V_K\otimes_k G)[-m,m])
       \subset C_n(G[a-m,b+m]).
\]

This proves the corrected inference at 1836–\allowbreak{}1840 with its original n,\allowbreak{}m,\allowbreak{}K,\allowbreak{}E,\allowbreak{}G,\allowbreak{} and completes the strong-generator argument.

For the next lemma,\allowbreak{} replace K by a coherent sheaf F. Smooth X is regular of finite dimension,\allowbreak{} so F is perfect locally; moreover F tensor\^{}L E\allowbreak{}=F tensor E because E is locally free. This is a coherent sheaf in degree zero. The bounds a\allowbreak{}=0,\allowbreak{}b\allowbreak{}=dim(X)\allowbreak{} apply to its cohomology,\allowbreak{} by the Noetherian dimension vanishing theorem (cohomology.tex,\allowbreak{} proposition-vanishing-Noetherian,\allowbreak{} current statement 4083–\allowbreak{}4086)\allowbreak{}. Thus the source\textquotesingle s m\allowbreak{}=max(m',\allowbreak{}m'\allowbreak{}+dim(X)\allowbreak{})\allowbreak{} still works. The corrected arbitrary-window assertion at 1862 must name H\textasciicircum{}i(X,\allowbreak{}F tensor E)\allowbreak{},\allowbreak{} with E the fixed sheaf from the preceding proof. Neither lemma\textquotesingle s final conclusion or dimension constant needs weakening.

The missing tensor factor is demonstrably substantive. On X\allowbreak{}=P\textasciicircum{}1\_\allowbreak{}k use homogeneous coordinates [X0:\allowbreak{}X1]\allowbreak{},\allowbreak{} t\allowbreak{}=X1/\allowbreak{}X0 on U0,\allowbreak{} and s\allowbreak{}=X0/\allowbreak{}X1\allowbreak{}=t\textasciicircum{}\{-1\} on U1. The frames X0\^{}e and X1\^{}e of O(e)\allowbreak{} obey X1\textasciicircum{}e\allowbreak{}=t\textasciicircum{}e X0\^{}e. The two-affine-cover Cech complex is

\[
 k[t]\oplus k[t^{-1}]
 \xrightarrow{(f,g)\mapsto t^e g-f}
 k[t,t^{-1}].
\]

All opens and their intersection are affine,\allowbreak{} so quasi-coherent affine acyclicity identifies its kernel and cokernel with H\^{}0 and H\textasciicircum{}1,\allowbreak{} with higher groups zero. For e\allowbreak{}=-1 the two sets of exponents are >\allowbreak{}=0 and <\allowbreak{}=-1. They are disjoint and their union is Z,\allowbreak{} so both kernel and cokernel vanish. For e\allowbreak{}=-2 their union omits exactly -1; the kernel is zero and the cokernel is k times t\^{}\{-1\}. Therefore

\[
 H^*(X,O(-1))=0,\qquad H^1(X,O(-2))=k.
\]

Take K\allowbreak{}=O(-1)\allowbreak{} and E\allowbreak{}=G\allowbreak{}=O\allowbreak{}+O(-1)\allowbreak{}. These E,\allowbreak{}G meet the actual diagonal condition:\allowbreak{} the determinant X0Y1-X1Y0 defines the diagonal and gives the exact resolution

\[
 0\to O(-1,-1)\xrightarrow{X_0Y_1-X_1Y_0}
 O\to O_\Delta\to0.
\]

On the coordinate charts the diagonal equation is the coordinate difference,\allowbreak{} or its invertible-coordinate counterpart,\allowbreak{} a non-zero-divisor; its quotient is precisely the diagonal chart. This proves injectivity,\allowbreak{} identifies the quotient and establishes the sequence globally. Both vector bundles are summands of E boxtimes G. The associated triangle O \allowbreak{}-> O\_\allowbreak{}Delta \allowbreak{}-> O(-1,\allowbreak{}-1)\allowbreak{}[1]\allowbreak{} puts O\_\allowbreak{}Delta in C\_\allowbreak{}2((E boxtimes G)\allowbreak{}[-1,\allowbreak{}1]\allowbreak{})\allowbreak{}.

Nevertheless K tensor E\allowbreak{}=O(-1)\allowbreak{}\allowbreak{}+O(-2)\allowbreak{} has nonzero H\textasciicircum{}1,\allowbreak{} while K has zero global cohomology in every degree. For example [0,\allowbreak{}0]\allowbreak{} bounds H\textasciicircum{}*(X,\allowbreak{}K)\allowbreak{} but does not bound V\_\allowbreak{}K. The stated inference from that range is therefore invalid.

The unqualified ``for any a<\allowbreak{}=b” claim in the next proof also fails. For any fixed m',\allowbreak{} choose F\allowbreak{}=O(-1)\allowbreak{} and a\allowbreak{}=b\allowbreak{}=m'\allowbreak{}+1. Its global cohomology is zero and satisfies that printed condition. Yet G is a sheaf in degree zero and the alleged generating window becomes [1,\allowbreak{}2m'\allowbreak{}+1]\allowbreak{}. Section 4 shows that every object of C\_\allowbreak{}n(G[1,\allowbreak{}2m'\allowbreak{}+1]\allowbreak{})\allowbreak{} has zero cohomology sheaf in degree zero. F has nonzero degree-zero cohomology sheaf O(-1)\allowbreak{},\allowbreak{} a contradiction. Using F tensor E in the premise excludes this false range and proves the intended assertion.

These two loci are recorded as one propagated missing-factor finding,\allowbreak{} with two exact operations. The 101 equiv.tex report routing records have no unlocated record and none overlaps 1836–\allowbreak{}1838 or 1862–\allowbreak{}1863. This bounded inventory check is not a claim to have reviewed the remaining Equivalences reports.

\subsection{\texorpdfstring{9. Propagation,\allowbreak{} scope and remaining work}{9.  scope and remaining work}}\label{proofs-10691-11108-propagation-scope-and-remaining-work}

The downstream boundedness lemma equiv.tex 1871–\allowbreak{}1888 remains valid. Choose the uniform m,\allowbreak{}n,\allowbreak{}G from the repaired diagonal lemma. For its homological functor choose the finite [a,\allowbreak{}b]\allowbreak{} supporting H\textasciicircum{}*(G)\allowbreak{},\allowbreak{} as allowed by that lemma\textquotesingle s assumptions. Section 4 gives the uniform range [a-m,\allowbreak{}b\allowbreak{}+m]\allowbreak{} for every coherent F. Enlarge it to [-M,\allowbreak{}M]\allowbreak{} with M\allowbreak{}=max(1,\allowbreak{}m-a,\allowbreak{}m\allowbreak{}+b)\allowbreak{}. This explicit M does not depend on the extension count n. For a cohomological functor apply the same reasoning with its reversed exact Hom sequence and its corresponding shift indices; a symmetric interval is unchanged by reversing signs.

The call to the classical-to-strong-generator lemma in more-algebra.tex 21561–\allowbreak{}21578 is preserved:\allowbreak{} if the cited earlier ring-specific lemmas give a classical generator R and any strong generator,\allowbreak{} Section 5 supplies the precise finite rn comparison. Only this categorical implication is checked here; this pass does not certify the unreviewed ring-specific premises.

All eleven physical reports in this span concern seven previously corrected units 0888–\allowbreak{}0894. P02-E0086 reports the unbound base-case index,\allowbreak{} although its range also touches the different truncation repair; its disposition is joined to 0889 by its actual text. P02-E0085 has the correct replacement but the inaccurate functor rationale discussed in Section 5. Neither range overlap nor an earlier admission is substituted for these semantic checks.

New source operations are limited to the shift-isomorphism qualification and the two missing-tensor-factor replacements. The sparse costs and bounded-window generator transfer remain separate editorial results with complete derivations above. This pass completes semantic review through 11108 and leaves Compact objects,\allowbreak{} starting 11121,\allowbreak{} next. Reading through 11170 is only read-ahead. No whole-chapter Equivalences or More on Algebra review,\allowbreak{} live-source composition,\allowbreak{} build or publication is claimed.
